\documentclass[11pt,a4paper]{article}

% Packages
\usepackage[utf8]{inputenc}
\usepackage[T1]{fontenc}
\usepackage{lmodern}
\usepackage{geometry}
\geometry{margin=2.5cm}
\usepackage{graphicx}
\usepackage{booktabs}
\usepackage{amsmath}
\usepackage{amssymb}
\usepackage{hyperref}
\usepackage{caption}
\usepackage{subcaption}
\usepackage{xcolor}
\usepackage{listings}
\usepackage{float}
\usepackage{enumitem}
\usepackage{microtype}

% Hyperref setup
\hypersetup{
    colorlinks=true,
    linkcolor=blue,
    citecolor=blue,
    urlcolor=blue,
    pdftitle={GPS Spoofing Detection Research Analysis},
    pdfauthor={Nikhil Sarwara}
}

% Code listings style
\lstset{
    basicstyle=\ttfamily\small,
    breaklines=true,
    frame=single,
    numbers=left,
    numberstyle=\tiny\color{gray},
    keywordstyle=\color{blue},
    commentstyle=\color{green!50!black},
    stringstyle=\color{orange},
    showstringspaces=false
}

% Custom commands
\newcommand{\todo}[1]{\textcolor{red}{\textbf{[TODO: #1]}}}

\title{\textbf{Research Analysis Draft}\\[0.5em]
\large GPS Spoofing Detection for PX4-Based Unmanned Aerial Vehicles}
\author{Nikhil Sarwara}
\date{SIT724 / SIT746 Research Project\\May 2026}

\begin{document}

\maketitle

%==============================================================================
\section{Introduction to Analysis}
%==============================================================================

The proliferation of unmanned aerial vehicles (UAVs) in civilian and military domains has elevated concerns regarding their vulnerability to Global Positioning System (GPS) spoofing attacks. In such attacks, adversaries transmit counterfeit GPS signals to manipulate a drone's perceived location, potentially leading to hijacking, denial of service, or safety hazards. Existing detection methods often rely on specialised hardware or single-point detection mechanisms, limiting their practicality for widespread deployment.

This research addresses the critical challenge of detecting GPS spoofing attacks on PX4-based drone systems through a software-only solution that leverages existing sensor telemetry. The purpose of this analysis is to evaluate whether the implemented detection system can reliably identify spoofing behaviour using data captured from simulated flight environments.

\textbf{Research Question.} \textit{Can a software-based system using standard PX4 telemetry and cross-sensor consistency checks detect GPS spoofing attacks in real time?}

To evaluate this question, the following hypotheses guide the analysis:
\begin{enumerate}[label=H\arabic*:,leftmargin=2em]
    \item GPS spoofing introduces measurable inconsistencies between GPS-reported position and inertial measurement unit (IMU) data.
    \item Heuristic-based labelling can generate sufficiently accurate training data without manual annotation.
    \item Machine learning models (Random Forest and 1D CNN) can distinguish normal flight windows from spoofed windows.
    \item The live inference pipeline can operate with latency suitable for real-time detection.
\end{enumerate}

This analysis focuses on telemetry capture, automated heuristic labelling, model evaluation, live inference behaviour, and the operational limitations of the system.

%==============================================================================
\section{Data Capture}
%==============================================================================

\subsection{Data Sources}

The primary data source for this analysis is telemetry generated by the PX4 Software-In-The-Loop (SITL) simulator coupled with the Gazebo physics engine. The simulation environment replicates realistic flight dynamics and sensor behaviour while allowing controlled injection of spoofing conditions. Table~\ref{tab:data_sources} summarises the data sources used.

\begin{table}[H]
\centering
\caption{Summary of data sources used in the analysis.}
\label{tab:data_sources}
\begin{tabular}{@{}lll@{}}
\toprule
\textbf{Source} & \textbf{Format} & \textbf{Purpose} \\
\midrule
PX4 SITL + Gazebo & MAVLink messages & Flight simulation and spoofing injection \\
GPS Monitor & Raw CSV logs & Telemetry collection at $\sim$10\,Hz \\
ML Pipeline & Processed CSV & Cleaned, labelled, windowed data \\
Live Inference & Timestamped CSV & Real-time detection outputs \\
\bottomrule
\end{tabular}
\end{table}

\subsection{Capture Methods and Procedures}

Data capture follows a structured pipeline. First, the PX4 SITL environment is launched within a Docker container, ensuring reproducible simulation conditions. The Gazebo simulator provides the physics backend, while QGroundControl or scripted MAVLink commands initiate flight missions. During flight, the \textbf{GPS Monitor} component connects to the PX4 autopilot via a MAVLink UDP bridge and records telemetry messages at approximately \textbf{10\,Hz}.

The captured telemetry includes GPS position (latitude, longitude, altitude), velocity vectors, acceleration, heading, and IMU angular rates. Data is stored as timestamped CSV files in the \texttt{gps\_logs/} directory. To create labelled attack scenarios, spoofing conditions are introduced in simulation by manipulating the GPS position stream while the IMU continues to report true vehicle motion.

\subsection{Rationale for Simulation-Based Capture}

Simulation-based data capture was chosen for several reasons:
\begin{itemize}
    \item \textbf{Safety and repeatability.} Controlled spoofing experiments pose no physical risk and can be repeated identically.
    \item \textbf{Scalability.} Large-volume datasets can be generated efficiently without flight-hour constraints.
    \item \textbf{Consistency.} The MAVLink protocol provides a standardised telemetry format across all runs.
    \item \textbf{Accessibility.} The same telemetry channels (GPS, IMU, velocity) are available on physical PX4 hardware, enabling future transferability.
\end{itemize}

\subsection{Reliability and Validity Measures}

Several measures were taken to enhance data quality:
\begin{itemize}
    \item \textbf{Controlled conditions.} Simulation parameters (wind, GPS noise, vehicle model) were held constant across runs to reduce uncontrolled variance.
    \item \textbf{Repeated runs.} Multiple flight missions were executed under both normal and spoofed conditions to ensure consistency.
    \item \textbf{Automated logging.} The GPS Monitor writes telemetry directly to CSV without manual transcription, eliminating human recording error.
    \item \textbf{Data cleaning.} Missing or malformed MAVLink messages are filtered during preprocessing.
\end{itemize}

\subsection{Limitations of Data Capture}

The following limitations are acknowledged:
\begin{itemize}
    \item Simulation data may not fully replicate real-world RF signal complexity, multipath effects, or sensor hardware noise.
    \item The current dataset contains a limited variety of spoofing scenarios (e.g., abrupt jumps vs. gradual drift).
    \item Label quality depends on the accuracy of heuristic rules, which may misclassify edge cases.
    \item The analysis may underrepresent borderline attack cases where spoofing intensity is low.
\end{itemize}

%==============================================================================
\section{Data Preparation}
%==============================================================================

\subsection{Preprocessing Pipeline}

Raw telemetry CSVs undergo a multi-stage preprocessing pipeline before model training. Figure~\ref{fig:pipeline} illustrates the workflow.

\begin{figure}[H]
\centering
\fbox{\parbox{0.9\textwidth}{\centering
\textbf{Raw Telemetry} $\rightarrow$ \textbf{Clean} $\rightarrow$ \textbf{Label} $\rightarrow$ \textbf{Window} $\rightarrow$ \textbf{Train / Val / Test}\\[0.3em]
\small
Raw CSV \quad $\rightarrow$ \quad Remove nulls, deduplicate \quad $\rightarrow$ \quad 8 heuristic rules \quad $\rightarrow$ \quad 30-message sliding windows \quad $\rightarrow$ \quad Stratified split
}}
\caption{Data preprocessing pipeline from raw telemetry to model input.}
\label{fig:pipeline}
\end{figure}

\textbf{Step 1 -- Cleaning.} Raw CSV files are inspected for missing values, duplicate timestamps, and out-of-range sensor readings. Invalid rows are removed or interpolated where appropriate.

\textbf{Step 2 -- Feature Extraction.} The following features are extracted or computed from the telemetry stream:
\begin{itemize}
    \item GPS position change (delta latitude, delta longitude, delta altitude)
    \item GPS velocity magnitude and direction
    \item Acceleration magnitude and change
    \item IMU angular rate stability
    \item Heading consistency between GPS track and IMU yaw
    \item GPS--IMU consistency score (physics-based heuristic)
\end{itemize}

\textbf{Step 3 -- Heuristic Labelling.} Eight physics-based heuristic rules are applied to label each telemetry record as \texttt{normal} or \texttt{spoofed}. The rules include GPS--IMU inconsistency checks, velocity anomalies, abrupt position jumps, and heading mismatches. A record is labelled \texttt{spoofed} if any heuristic flags an anomaly.

\textbf{Step 4 -- Windowing.} Temporal sliding windows of \textbf{30 messages} ($\sim$3\,seconds at 10\,Hz) are created to capture sequential patterns. Each window inherits the majority label of its constituent records.

\textbf{Step 5 -- Splitting.} Windowed data is divided into training, validation, and test sets. \todo{Add exact split ratios and counts. Ensure no data leakage by splitting on flight runs, not individual windows.}

\subsection{Avoiding Data Leakage}

To prevent information leakage, windows originating from the same flight run are kept together in either the training or test set, never split across both. \todo{Confirm this was implemented; if not, state it as a limitation and plan to fix.}

%==============================================================================
\section{Data Analysis Methods}
%==============================================================================

\subsection{Descriptive Statistics}

The dataset is characterised by descriptive statistics including total record counts, class distribution (normal vs. spoofed), sampling rate consistency, and flight duration distributions.

\subsection{Classification Evaluation}

Model performance is evaluated using standard binary classification metrics:
\begin{itemize}
    \item \textbf{Accuracy:} proportion of correctly classified windows.
    \item \textbf{Precision:} proportion of predicted spoofed windows that are truly spoofed.
    \item \textbf{Recall:} proportion of actual spoofed windows correctly detected.
    \item \textbf{F1-Score:} harmonic mean of precision and recall.
    \item \textbf{Confusion Matrix:} true/false positive and negative counts.
\end{itemize}

These metrics were chosen because accuracy alone is insufficient for safety-critical applications where false negatives (missed spoofing) and false positives (unnecessary alerts) carry different costs.

\subsection{Model Comparison}

Two models are compared:
\begin{enumerate}
    \item \textbf{Random Forest (RF):} An ensemble of decision trees offering fast inference and interpretable feature importance.
    \item \textbf{1D Convolutional Neural Network (CNN):} A deep learning model capturing temporal patterns within each 30-message window.
\end{enumerate}

\subsection{Live Inference Analysis}

Live inference logs from \texttt{ui/logs/} are analysed to assess:
\begin{itemize}
    \item Detection latency between spoofing onset and model alert.
    \item Stability of predictions across consecutive windows.
    \item Frequency of false positives during normal flight segments.
\end{itemize}

\subsection{Breaking-Point Analysis (If Available)}

If spoofing intensity was varied systematically, a breaking-point analysis characterises the attack severity threshold at which detection performance degrades. \todo{Include if data available; otherwise state as planned future work.}

%==============================================================================
\section{Results}
%==============================================================================

\subsection{Dataset Summary}

\todo{Fill in exact numbers from project logs.}

\begin{table}[H]
\centering
\caption{Dataset summary.}
\label{tab:dataset}
\begin{tabular}{@{}lr@{}}
\toprule
\textbf{Metric} & \textbf{Value} \\
\midrule
Total flight runs & \todo{Add count} \\
Total telemetry records & \todo{Add count} \\
Sampling rate & 10\,Hz \\
Normal records & \todo{Add count} \\
Spoofed records & \todo{Add count} \\
Total windows generated & \todo{Add count} \\
Window size & 30 messages ($\sim$3\,s) \\
Training windows & \todo{Add count} \\
Validation windows & \todo{Add count} \\
Test windows & \todo{Add count} \\
\bottomrule
\end{tabular}
\end{table}

\subsection{Label Generation Results}

Figure~\ref{fig:labels} shows an example telemetry segment with heuristic labels overlaid. The heuristic labels align with the injected spoofing periods, confirming that the physics-based rules capture the intended attack behaviour.

\begin{figure}[H]
\centering
\fbox{\parbox{0.8\textwidth}{\centering
\vspace{3em}
\textit{[Insert \texttt{labels vs time} figure from project figures]}
\vspace{3em}
}}
\caption{Heuristic labels vs. time for a representative flight segment.}
\label{fig:labels}
\end{figure}

\subsection{Model Performance}

Table~\ref{tab:model_perf} presents the classification results for both models.

\begin{table}[H]
\centering
\caption{Model performance comparison.}
\label{tab:model_perf}
\begin{tabular}{@{}lcccc@{}}
\toprule
\textbf{Model} & \textbf{Accuracy} & \textbf{Precision} & \textbf{Recall} & \textbf{F1} \\
\midrule
Random Forest & \todo{Add} & \todo{Add} & \todo{Add} & \todo{Add} \\
1D CNN & \todo{Add} & \todo{Add} & \todo{Add} & \todo{Add} \\
\bottomrule
\end{tabular}
\end{table}

\todo{Insert confusion matrix figures for both models.}

\subsection{Live Inference Results}

Recent live inference logs (e.g., \texttt{live\_20260510\_050706.csv}) demonstrate the detector operating on streaming telemetry. \todo{Add one example timeline figure showing detection timestamps relative to spoofing onset.}

\subsection{System Behaviour Under Attack}

During spoofed flight segments, the GPS position stream deviated from the IMU-reported motion, producing a measurable inconsistency score. The detector flagged these windows as anomalous, while normal flight segments with aggressive manoeuvres (where both GPS and IMU agreed) were correctly classified as normal. \todo{Add specific example or quantitative summary.}

\subsection{Breaking Point / Sensitivity}

\todo{If available, present the attack intensity threshold where detection degrades. Otherwise state: ``Breaking-point analysis is planned as future work to quantify the minimum spoofing intensity detectable by the system.''}

%==============================================================================
\section{Discussion and Interpretation}
%==============================================================================

\subsection{Interpretation of Findings}

The results support the feasibility of a software-only GPS spoofing detection framework for PX4 drones. The telemetry pipeline successfully captured sufficient data to identify spoofing behaviour in controlled simulation conditions. The heuristic labelling pipeline reduced reliance on manual annotation, enabling scalable dataset generation.

The Random Forest model offers fast, interpretable detection suitable for resource-constrained embedded systems, while the 1D CNN captures temporal dependencies that may improve robustness against gradual spoofing attacks. The GPS--IMU consistency heuristic is particularly valuable because it distinguishes spoofing from natural disturbances (e.g., wind gusts) conceptually: wind affects both GPS and IMU, whereas spoofing affects GPS alone.

\subsection{Careful Interpretation of High Accuracy}

\todo{If 100\% test accuracy is reported, include the following paragraph. Otherwise adjust.}

The models achieved high test accuracy, suggesting strong separability between normal and spoofed windows in the current dataset. However, this result must be interpreted cautiously. Controlled simulation conditions, limited scenario diversity, and potential data leakage risks (if windows from the same run were split across sets) may inflate performance metrics. Therefore, while the result is promising, it is not yet sufficient to claim general real-world robustness. Future work must validate performance on more diverse datasets and physical hardware.

\subsection{Unexpected Findings}

Several observations warrant further investigation:
\begin{itemize}
    \item \textbf{Gradual spoofing.} Abrupt position jumps were detected reliably, but gradual drift scenarios may be harder to distinguish from natural sensor noise.
    \item \textbf{Heuristic dominance.} Certain heuristics (e.g., GPS--IMU inconsistency) may dominate the labelling, potentially biasing the training data.
    \item \textbf{Live inference lag.} The 30-message window introduces a detection delay of approximately 3\,seconds, which may be acceptable for some missions but critical for high-speed operations.
    \item \textbf{False positives during manoeuvres.} Aggressive turns briefly decouple GPS and IMU readings, which could trigger false alerts if the consistency threshold is too sensitive.
\end{itemize}

%==============================================================================
\section{Limitations and Threats to Validity}
%==============================================================================

The following limitations affect the generalisability and validity of the findings:

\begin{enumerate}
    \item \textbf{Simulation-only data.} The analysis is based entirely on PX4 SITL telemetry. Real-world RF environments, sensor hardware noise, and multipath effects are not represented.
    \item \textbf{Limited spoofing diversity.} The current dataset primarily covers abrupt spoofing scenarios. Gradual drift, intermittent spoofing, and multi-sensor attacks are underrepresented.
    \item \textbf{Heuristic label bias.} The eight heuristic rules may misclassify edge cases or dominate labels unevenly, introducing training bias.
    \item \textbf{Incomplete IMU integration.} While the GPS--IMU consistency check is used for labelling, full IMU features are not yet integrated into the ML pipeline training input.
    \item \textbf{Potential overfitting.} High accuracy on a limited dataset may indicate overfitting rather than generalisable performance.
    \item \textbf{Latency measurement.} Formal real-time latency benchmarks have not yet been conducted; the 3-second window delay is a theoretical estimate.
    \item \textbf{Breaking point unknown.} The minimum spoofing intensity required for reliable detection has not been quantified.
    \item \textbf{No physical validation.} The system has not been tested on real PX4 hardware or in live RF conditions.
\end{enumerate}

%==============================================================================
\section{Conclusion}
%==============================================================================

This analysis demonstrates that the implemented PX4-based telemetry pipeline can capture sufficient data to identify GPS spoofing behaviour in simulation. The combination of heuristic labelling, temporal windowing, and machine learning classification produced strong separation between normal and spoofed flight conditions. These findings support the feasibility of a software-only spoofing detection framework, while also highlighting the need for broader datasets, stronger cross-sensor validation, and real-world testing.

Specifically:
\begin{itemize}
    \item The data capture methodology is repeatable, scalable, and uses standard PX4 telemetry channels.
    \item The heuristic labelling pipeline generates training data without manual annotation, though its accuracy depends on rule design.
    \item Both Random Forest and 1D CNN models classify spoofed windows with high accuracy in controlled conditions.
    \item The GPS--IMU consistency check provides a physics-based discriminator that conceptually separates spoofing from natural disturbances.
\end{itemize}

Future work will focus on integrating full IMU data into the ML pipeline, expanding dataset diversity, measuring formal latency benchmarks, quantifying the detection breaking point, and validating the system on physical PX4 hardware.

%==============================================================================
% APPENDIX (Optional)
%==============================================================================
\appendix
\section{Sample Log Format}
\label{app:logs}

\todo{Include a small excerpt of raw and processed CSV logs if helpful for markers.}

\section{Hyperparameters}
\label{app:hyperparams}

\todo{List model hyperparameters: RF trees, CNN layers, window size, learning rate, etc.}

\end{document}
